% ============================================================
%  Coherence Field Framework (CFF)
%  Formal Specification
%  Version 1.0 — May 2026
%  Confidential — Trade Secret — Do Not Distribute
%  Overleaf-compatible LaTeX
% ============================================================
\documentclass[11pt, a4paper]{article}

\usepackage[margin=2.2cm]{geometry}
\usepackage{amsmath, amssymb, amsthm, mathtools}
\usepackage{booktabs, array, longtable}
\usepackage{xcolor}
\usepackage{hyperref}
\usepackage{titlesec}
\usepackage{enumitem}
\usepackage{fancyhdr}
\usepackage{setspace}
\usepackage{listings}

\hypersetup{
  colorlinks=true,
  linkcolor=blue!50!black,
  urlcolor=blue!60!black,
}

\definecolor{nm-dark}{RGB}{30, 30, 30}
\definecolor{nm-gray}{RGB}{100, 100, 100}
\definecolor{nm-green}{RGB}{20, 100, 40}
\definecolor{nm-red}{RGB}{160, 30, 30}
\definecolor{nm-orange}{RGB}{180, 100, 20}
\definecolor{code-bg}{RGB}{245, 245, 245}
\definecolor{code-kw}{RGB}{0, 0, 180}
\definecolor{code-str}{RGB}{0, 120, 0}
\definecolor{code-cmt}{RGB}{120, 120, 120}

\pagestyle{fancy}\fancyhf{}
\rhead{\textcolor{nm-gray}{\small CFF Formal Specification v1.0}}
\lhead{\textcolor{nm-gray}{\small Trade Secret --- Do Not Distribute}}
\cfoot{\thepage}

\titleformat{\section}{\large\bfseries\color{nm-dark}}{\thesection}{0.6em}{}[\titlerule]
\titleformat{\subsection}{\normalsize\bfseries\color{nm-dark}}{\thesubsection}{0.6em}{}
\titleformat{\subsubsection}{\normalsize\itshape\color{nm-dark}}{\thesubsubsection}{0.6em}{}

\setlength{\parindent}{0pt}\setlength{\parskip}{6pt}\onehalfspacing

\lstset{
  language=Python,
  basicstyle=\ttfamily\small,
  backgroundcolor=\color{code-bg},
  keywordstyle=\color{code-kw}\bfseries,
  stringstyle=\color{code-str},
  commentstyle=\color{code-cmt}\itshape,
  breaklines=true, frame=single, framesep=4pt,
  xleftmargin=6pt, xrightmargin=6pt,
  numbers=none, showstringspaces=false, tabsize=4,
}

\newtheorem{definition}{Definition}[section]
\newtheorem{axiom}{Axiom}[section]
\newtheorem{theorem}{Theorem}[section]
\newtheorem{proposition}{Proposition}[section]
\newtheorem{corollary}{Corollary}[section]
\newtheorem{remark}{Remark}[section]
\newtheorem{prediction}{Prediction}[section]

\newcommand{\pass}{\textcolor{nm-green}{\textbf{PASS}}}
\newcommand{\parti}{\textcolor{nm-orange}{\textbf{PARTIAL}}}
\newcommand{\fail}{\textcolor{nm-red}{\textbf{FAIL}}}
\newcommand{\untested}{\textcolor{nm-orange}{\textbf{UNTESTED}}}

\begin{document}

\begin{center}
  {\LARGE \textbf{Coherence Field Framework}}\\[4pt]
  {\LARGE \textbf{(CFF)}}\\[8pt]
  {\large Formal Specification}\\[4pt]
  {\large From Theoretical Foundation to Applied Discovery}\\[12pt]
  {\color{nm-gray}\small Version 1.0 \quad|\quad May 2026}\\[4pt]
  {\color{nm-gray}\small Joseph Forrester \quad|\quad DSF-AI Architecture}\\[4pt]
  \textcolor{nm-red}{\textbf{TRADE SECRET --- CONFIDENTIAL --- DO NOT DISTRIBUTE}}
\end{center}

\vspace{1em}\hrule\vspace{1em}

\begin{center}
\fbox{\parbox{0.88\textwidth}{\centering
\textbf{Classification:} This document is a trade secret of Joseph Forrester /
DSF-AI. It contains the complete formal framework connecting the Universal Field
of Coherent Processes (UFCP) theoretical foundation to the Coherence Field
Framework (CFF) applied discovery methodology. The kernel implementation,
axiom set, structural filters, and forcing chain methodology are proprietary.
Nothing in this document may be reproduced, distributed, or disclosed without
written authorization.}}
\end{center}

\newpage
\tableofcontents
\newpage

% ============================================================
% PART 0: EXECUTIVE SUMMARY
% ============================================================
\section*{Executive Summary}
\addcontentsline{toc}{section}{Executive Summary}

The Coherence Field Framework (CFF) is the applied successor to the
Universal Field of Coherent Processes (UFCP). Where UFCP provides the
theoretical physics --- a conservative field framework deriving classical,
quantum, and relativistic dynamics from a single Hermitian action ---
CFF extracts from that theory a set of \textbf{structural axioms} and
converts them into \textbf{computable filters} for material and system
discovery.

\subsection*{What CFF Is}

CFF is a deterministic discovery methodology. Given a target physical
property and a materials database, CFF outputs a ranked list of
candidates with a fully auditable forcing chain. Every rejected
candidate has a traceable reason. Every surviving candidate has a
documented path through the filters.

\subsection*{What CFF Is Not}

CFF is not a simulation tool, a machine learning model, or a
curve-fitting engine. It does not predict numerical values from
first principles (that is UFCP's role). CFF identifies
\textit{architectural classes} --- the structural families within
which solutions must exist --- by sequential constraint elimination.

\subsection*{Relationship to UFCP}

\begin{center}
\begin{tabular}{@{}lll@{}}
\toprule
& \textbf{UFCP} & \textbf{CFF} \\
\midrule
\textbf{Role} & Theoretical foundation & Applied discovery tool \\
\textbf{Inputs} & Physical constants, action functional &
  Empirical databases, target property \\
\textbf{Outputs} & Equations of motion, predictions &
  Ranked candidates, forcing chains \\
\textbf{Method} & Variational calculus, Noether currents &
  Sequential constraint elimination \\
\textbf{Free parameters} & One ($\rho_0$) & Zero (thresholds
  from empirical data) \\
\textbf{Validation} & 111 tests, zero falsifications &
  2 target domains demonstrated \\
\bottomrule
\end{tabular}
\end{center}

\subsection*{Demonstrated Applications}

\begin{enumerate}[nosep]
  \item \textbf{Room-temperature superconductivity (RTSC):}
    CFF independently derived Cu$_3$O$_2$ on LaAlO$_3$(111) as
    the unique survivor of six structural filters applied to
    the entire known materials space. Patent filed April 2026
    (US Provisional 64/037,691).
  \item \textbf{Nanoscale heat transfer metamaterials:}
    CFF identified 12 all-phononic metamaterial geometries
    predicted to outperform the plasmonic-phononic approach
    published in \textit{Nature} (Wang et al., 2026).
  \item \textbf{Nanoparticle property prediction:}
    The UFCP geometric framework produced 63 quantitative
    predictions across 7 elements and 4 property types with
    6.5\% average error and zero free parameters.
  \item \textbf{Phase transition detection:}
    The UF-Core computational kernel (L0--L4) detected phase
    transitions in 7 materials across 4 physical phenomena
    with zero parameter changes.
\end{enumerate}

% ============================================================
% PART I: THEORETICAL FOUNDATION
% ============================================================
\newpage
\part{Theoretical Foundation (UFCP)}

\section{Scope}

This part summarizes the UFCP theoretical framework from which
CFF axioms are derived. The complete mathematical treatment is
in the UFCP Comprehensive Specification v1.0 and Not-Math
Framework v2.0. This section extracts only what CFF requires.

% ============================================================
\section{The Coherence Field}

\begin{definition}[Coherence Field]
\label{def:coherence-field}
Let $\psi = \sqrt{\rho}\, e^{i\phi}$ be the coherence field on
a Lorentzian manifold $(M, g_{\mu\nu})$ with minimal gauge coupling
$D_\mu = \partial_\mu + i\frac{q}{\hbar}A_\mu$.

The field carries three physical attributes:
\begin{itemize}[nosep]
  \item $\rho(\mathbf{x}, t) \geq 0$: coherent-potential density
  \item $\phi(\mathbf{x}, t) \in \mathbb{R}$: phase
  \item $\omega = \partial_t \phi$: local tempo
\end{itemize}
\end{definition}

\begin{axiom}[UFCP Action]
\label{ax:action}
\begin{equation}
  \boxed{
  S = \int d^4x\, \sqrt{-g} \left[
    \frac{i\hbar}{2}\!\left(\psi^* D_t\psi - \psi D_t\psi^*\right)
    - \psi^* \hat{\Omega}_g \psi
    - U(\rho)
    - \frac{1}{2}\int \rho\, W\, \rho'\, \sqrt{-g'}\, d^4x'
  \right]
  }
\end{equation}
where $\hat{\Omega}_g$ is a Hermitian kinetic operator, $U(\rho)$
is the local self-interaction potential (bounded below), and
$W(\mathbf{x}-\mathbf{x}')$ is the symmetric interaction kernel.
\end{axiom}

\begin{axiom}[Kernel Symmetry]
\label{ax:kernel-sym}
$W(\mathbf{x}-\mathbf{x}') = W(\mathbf{x}'-\mathbf{x})$.
Grounded in CPT invariance.
\end{axiom}

\subsection{Equation of Motion}

Variation of $S$ with respect to $\psi^*$:
\begin{equation}
  \boxed{
  i\hbar D_t \psi = \hat{\Omega}_g \psi
    + \frac{\partial U}{\partial \rho}\,\psi
    + (W * \rho)\,\psi
  }
\end{equation}

\subsection{Conserved Currents}

Global $U(1)$ symmetry gives $\nabla_\mu J^\mu = 0$ with
$J^0 = \rho$ and $\mathbf{J} = \rho\,\mathbf{v}$.

\subsection{Correspondence Limits}

The UFCP action reproduces exactly:
\begin{itemize}[nosep]
  \item Schr\"odinger equation ($U = 0$, $W = 0$, non-relativistic
    $\hat{\Omega}_g$)
  \item Newtonian mechanics (Ehrenfest limit, decoherence via $W$)
  \item Relativistic wave dynamics (Klein--Gordon / Dirac
    $\hat{\Omega}_g$)
  \item Born's rule (from Gleason's theorem under four structural
    measurement constraints)
\end{itemize}

% ============================================================
\section{The Geometric Coupling Law}

\begin{theorem}[Geometric Coupling Law]
\label{thm:gcl}
\begin{equation}
  \boxed{Q = Q_0\left(1 + \alpha^2 \lambda^2 \Gamma\right)}
\end{equation}
where $Q_0$ is the uncoupled value, $\alpha = 1/137.036$ is
the fine structure constant (derived from 3D soliton stability:
$\alpha = 1/N_c^2$), $\lambda$ is the material-specific coupling
parameter, and $\Gamma = \cos\theta$ is the geometry factor.

Ring geometry ($\theta = 0$): $\Gamma = 1$ (maximum coupling).
Sphere geometry ($\theta = \pi/2$): $\Gamma = 0$ (no coupling
correction). Zero free parameters.
\end{theorem}

\textbf{Validation:} 5 independent physical systems, 24+
materials, all consistent. See UFCP Comprehensive Spec v1.0
Part~III.

% ============================================================
\section{The Soliton Ontology}

Particles in UFCP are solitonic density concentrations of
$\psi$, not fundamental point objects. The soliton width matches
the Compton wavelength when $\lambda\rho_0 = mc^2/2$.

\textbf{Key consequence for CFF:} Structural stability of a
material phase is a soliton stability problem. Phase transitions
correspond to soliton bifurcations. CFF's structural filters
are constraints on soliton stability conditions.

% ============================================================
\section{Connected Constants}

From the UFCP action, the following physical constants are
determined (not free):

\begin{center}
\begin{tabular}{@{}lll@{}}
\toprule
\textbf{Constant} & \textbf{Determined by} & \textbf{Status} \\
\midrule
$\alpha$ & 3D soliton stability ($1/N_c^2$) & Derived \\
$c$ & $G$, $\rho_0$, $\alpha$, $\hbar$ & Derived \\
$G$ & $\lambda/(4\pi\rho_0)$ & Derived \\
$\rho_{DE}$ & $m\rho_0/2$ & Derived \\
$\rho_0$ & Initial conditions & Free parameter \\
$\hbar$ & Coherence field quantum & Unit convention \\
\bottomrule
\end{tabular}
\end{center}

One free parameter. One unit convention. Everything else
determined.

% ============================================================
% PART II: CFF AXIOM SET
% ============================================================
\newpage
\part{CFF Axiom Set}

\section{Scope}

The CFF axioms are structural constraints extracted from the
UFCP action and its consequences. Each axiom generates one or
more computable filters for the discovery algorithm. The axioms
are numbered A1--A8.

\textbf{Design principle:} Each CFF axiom can be independently
justified on conventional physics grounds without requiring
acceptance of UFCP. CFF is designed to be auditable by domain
experts who may reject the theoretical foundation entirely.

% ============================================================
\section{Axiom Definitions}

\begin{axiom}[A1 --- Structural Existence]
\label{ax:a1}
A physical system must admit stable localized excitations
(solitonic states) in its effective dimensionality. Systems
whose dimensionality does not support topological closure of
solitonic modes cannot sustain coherent structural phases.

\textbf{Filter:} Reject systems with dimensionality that does
not support stable solitons ($d < 2$ for topological solitons;
$d = 1$ supported only with specific nonlinearities).

\textbf{Conventional justification:} Mermin--Wagner theorem
(no continuous symmetry breaking in $d \leq 2$ without
topological protection); BKT theory (topological defects
stabilize 2D systems).
\end{axiom}

\begin{axiom}[A2 --- Norm Conservation]
\label{ax:a2}
The total coherent-potential density $\int \rho\, d^dx$ is
conserved under time evolution. Systems that violate norm
conservation cannot maintain structural identity.

\textbf{Filter:} Reject candidate materials or frameworks that
require non-unitary evolution (e.g., open quantum systems treated
as fundamental rather than effective).

\textbf{Conventional justification:} Unitarity of quantum
mechanics; charge conservation; continuity equation.
\end{axiom}

\begin{axiom}[A3 --- Boundary Condition Compatibility]
\label{ax:a3}
The boundary conditions of a structural system must not impose
competing winding structures that destructively interfere with
the target phase. Substrate--film interfaces must be
thermodynamically inert with respect to the structural phase.

\textbf{Filter:} Reject substrates that (i)~react with the
film material, (ii)~impose a competing electronic or magnetic
phase, or (iii)~have lattice mismatch exceeding the coherence
length of the target phase.

\textbf{Conventional justification:} Epitaxial growth requires
lattice matching; competing phases destroy long-range order;
interfacial reactions create dead layers.
\end{axiom}

\begin{axiom}[A4 --- Deterministic Evolution]
\label{ax:a4}
The structural state of a system evolves deterministically from
initial conditions. The perception of structure must be
reproducible: the same input field must produce the same
structural description.

\textbf{Filter:} Reject stochastic methods (Monte Carlo,
random initialization) as structural analysis tools. Accept
only deterministic, auditable algorithms.

\textbf{Conventional justification:} Reproducibility is
a foundational requirement of the scientific method.
\end{axiom}

\begin{axiom}[A5 --- Dispositional Preservation]
\label{ax:a5}
Carrier tuning (doping, gating, pressure) must not disorder
the coherence lattice. The structural geometry that produces
the target property must be preserved under the tuning method.

\textbf{Filter:} Reject chemical doping that introduces
substitutional disorder. Accept electrostatic gating, strain
tuning, and interface doping that modify carrier density
without disrupting the lattice.

\textbf{Conventional justification:} Chemical doping in
cuprates introduces pair-breaking scattering; ionic liquid
gating preserves lattice order (measured by RHEED, XRD).
\end{axiom}

\begin{axiom}[A6 --- Non-Trivial Winding]
\label{ax:a6}
The structural geometry must support non-trivial phase winding
--- topological defects that cannot be continuously deformed to
trivial configurations. Geometrically frustrated lattices
satisfy this axiom; non-frustrated lattices do not.

\textbf{Filter:} Reject square, cubic, BCC, and other
non-frustrated geometries. Accept kagome, triangular,
honeycomb, pyrochlore, and hyperkagome geometries.

\textbf{Conventional justification:} Geometric frustration
suppresses N\'eel ordering and enables spin-liquid,
superconducting, and topological phases. Kagome lattice
superconductors (CsV$_3$Sb$_5$, Cu-BHT) have weak-coupling
ratios; square-lattice cuprates have strong-coupling ratios.
\end{axiom}

\begin{axiom}[A7 --- Coherence Preservation]
\label{ax:a7}
The coupling regime must preserve topological winding as the
conserved quantity, not pair binding energy. Strong coupling
destroys topological coherence by localizing pairs.

\textbf{Filter:} Reject systems with BCS coupling ratio
$2\Delta/(k_B T_c) > r_{max}$ (where $r_{max}$ is
target-dependent; $r_{max} = 5.0$ for RTSC, $6.0$ for HTSC).

\textbf{Conventional justification:} Known weak-coupling
superconductors (MgB$_2$: ratio 3.5; CsV$_3$Sb$_5$: ratio
$\sim$4.0) achieve higher $T_c$-to-gap efficiency than
strong-coupling systems (YBCO: ratio 5.6--6.0; Bi-2212:
ratio 5.5--7.0).
\end{axiom}

\begin{axiom}[A8 --- Local Field Response]
\label{ax:a8}
The coupling amplitude between neighboring coherent sites
must exceed a threshold set by the target property. Below
this threshold, structural coherence cannot propagate across
the lattice at the required energy scale.

\textbf{Filter:} Reject systems with superexchange $J$ (or
equivalent coupling constant) below the target-specific
threshold ($J \geq 100$~meV for RTSC, $J \geq 50$~meV for
HTSC, phonon Q $\times$ oscillator strength $\geq$ threshold
for thermal metamaterials).

\textbf{Conventional justification:} Superexchange $J$
determines the pairing energy scale. Published values of $J$
for all transition-metal oxide bridges are tabulated (YBCO
Cu-O-Cu: 130~meV; NiO Ni-O-Ni: 85~meV; herbertsmithite
Cu-OH-Cu: 170~meV). Only Cu-based bridges exceed 100~meV.
\end{axiom}

% ============================================================
\section{Axiom Summary}

\begin{center}
\begin{longtable}{@{}clp{5cm}p{4cm}@{}}
\toprule
\textbf{Axiom} & \textbf{Name} & \textbf{CFF Filter} &
  \textbf{UFCP Origin} \\
\midrule
\endhead
A1 & Structural Existence & Dimensionality must support
  stable solitons & Soliton stability in $d$ dimensions \\
A2 & Norm Conservation & Unitary evolution required &
  $U(1)$ Noether current \\
A3 & Boundary Compatibility & Substrate must be inert and
  lattice-matched & Boundary conditions on $\psi$ \\
A4 & Deterministic Evolution & Reproducible structural
  perception & Deterministic EoM from $\delta S = 0$ \\
A5 & Dispositional Preservation & Doping must not disorder
  lattice & $J$ preserved under $W$ perturbation \\
A6 & Non-Trivial Winding & Frustrated geometry required &
  Topological soliton stability \\
A7 & Coherence Preservation & Weak coupling regime &
  Winding number as conserved quantity \\
A8 & Local Field Response & Coupling amplitude above
  threshold & $W * \rho$ exceeds $U(\rho)$ barrier \\
\bottomrule
\end{longtable}
\end{center}

% ============================================================
% PART III: COMPUTATIONAL KERNEL (L0-L4)
% ============================================================
\newpage
\part{Computational Kernel (L0--L4)}

\section{Scope}

The UF-Core Kernel is the computational implementation of CFF's
structural perception. It takes any dimensionless ordered field
and outputs a description of that field's structural geometry.
The kernel is domain-agnostic: it does not know what the field
represents.

\section{Layer Architecture}

\begin{center}
\begin{tabular}{@{}clp{7cm}@{}}
\toprule
\textbf{Layer} & \textbf{Name} & \textbf{Function} \\
\midrule
L0 & Structural Field Normalization &
  Converts raw input to Structural Event Vector (SEV):
  $(F, \Delta F, \sigma, \kappa)$. Normalized field value,
  first derivative, structural change rate, curvature. \\
L1 & Gate Formation \& Segmentation &
  Segments the field into \textit{gates} --- contiguous
  intervals of internally consistent structural character.
  Gate boundary operator:
  $D(t) = \alpha_1\|\Delta F\| + \alpha_2\sigma + \alpha_3\kappa$.
  New gate opens when $D(t) < \tau_D$. \\
L2 & Interpretive Structural Field &
  Maps gate sequences to interpretive vectors. Classifies
  each gate as STABLE, VOLATILE, TRANSITIONAL, or DEGENERATE.
  Performs mosaic projection. \\
L3 & Resonance Functional &
  Computes resonance vector:
  $R^{\text{vec}}_k = (w_k, \|CV_k\|/\|CV\|_{\max}, S_k,
  1/(1+C_k), 1-U_k)$.
  Unified Resonance Field: $URF_k = g_k \cdot R(k)$.
  Contextual gating with hysteresis. \\
L4 & Decision Dynamics &
  Outputs Decision Surface Field (DSF), a 7-tuple per gate:
  $(D_k, M_k, U^*_k, B_k, P_k, C_k, R_{\text{rev},k})$. \\
\bottomrule
\end{tabular}
\end{center}

\section{DSF 7-Tuple (L4 Output)}

\begin{definition}[Decision Surface Field]
\label{def:dsf}
The DSF is the kernel's sole output. Each gate $k$ receives
a 7-tuple:

\begin{center}
\begin{tabular}{@{}clp{6cm}@{}}
\toprule
\textbf{Symbol} & \textbf{Name} & \textbf{Meaning} \\
\midrule
$D_k$ & Direction & Structural field direction
  ($+1$, $0$, $-1$). Reversal marks a structural
  direction change. \\
$M_k$ & Momentum & Rate of structural change. How fast
  the field geometry is evolving. \\
$U^*_k$ & Uncertainty & Confidence in structural
  assessment. Peaks at $\sim$0.667 at phase transitions
  (maximum structural ambiguity). \\
$B_k$ & Breathing & Structural freedom / expansion.
  How much room the field has to move. \\
$P_k$ & Pressure & Transition sensitivity. Spikes when
  structural state is under stress. \\
$C_k$ & Complexity & Mosaic divergence. How many distinct
  structural projections exist within a gate. \\
$R_{\text{rev},k}$ & Reversal & How often the structural
  direction flips within a gate. High = unstable. \\
\bottomrule
\end{tabular}
\end{center}
\end{definition}

\section{Domain-Agnostic Property}

The same L0--L4 kernel, with identical parameters, has been
applied to:
\begin{itemize}[nosep]
  \item Financial time series (TFE: 11,884 symbols, 5 years)
  \item Superconductor R(T) curves (7 materials, 4 phenomena)
  \item Pharmaceutical DSC data (polymorph screening)
  \item Neural network weight distributions (MACE audit)
  \item Nanoparticle property data (cluster physics)
  \item IR sensor calibration data (ArcLoom robotics)
\end{itemize}

No parameter changes between domains. The kernel perceives
structure, not content.

\section{Universal Uncertainty Ceiling}

In every material tested, the kernel uncertainty $U^*_k$
peaked at $\mathbf{0.667}$ at the phase transition point.
This value was not programmed --- it emerges from the structural
geometry of the kernel when the input field is maximally
ambiguous. This universality across different physical phenomena
is consistent with the kernel measuring a fundamental property
of structural phase transitions.

% ============================================================
% PART IV: DISCOVERY ALGORITHM
% ============================================================
\newpage
\part{Discovery Algorithm}

\section{Scope}

The CFF Discovery Algorithm converts the axiom set (Part~II)
into a deterministic, auditable material discovery tool. Given
a target physical property and empirical databases, the algorithm
outputs ranked candidates with forcing chains.

\section{Four-Stage Pipeline}

\begin{enumerate}
  \item \textbf{Stage 1: CFF Structural Filters.}
    Apply axioms A1--A8 sequentially. Each axiom generates a
    filter that rejects candidates violating that axiom. Coarser
    filters first. Each rejected candidate is logged with the
    specific axiom and threshold that eliminated it.

  \item \textbf{Stage 2: Class Definition.}
    The surviving filter constraints define an
    \textit{architectural class} --- a structural family within
    which solutions must exist. This is the algorithm's primary
    output.

  \item \textbf{Stage 3: Class Enumeration.}
    Given the class definition, enumerate candidate materials
    from chemistry databases. Cross-reference with synthesis
    precedents from published literature. Rank by precedent
    strength.

  \item \textbf{Stage 4: Output.}
    Produce ranked candidates with forcing chains, rejected
    families with reasons, and the architectural class
    definition. All outputs are deterministic and reproducible.
\end{enumerate}

\section{Required Databases}

\begin{center}
\begin{tabular}{@{}lp{7cm}l@{}}
\toprule
\textbf{Database} & \textbf{Contents} & \textbf{Source} \\
\midrule
Bridge & Metal-ligand-metal superexchange constants $J$ &
  Published measurements \\
Substrate & Lattice parameters, (111) surface mesh,
  thermodynamic inertness & Crystallography databases \\
Composition & Demonstrated or proposed target-geometry
  compositions & Published synthesis papers \\
Target & BCS ratio thresholds, $J$ thresholds, domain-specific
  figures of merit & Empirical condensed matter \\
\bottomrule
\end{tabular}
\end{center}

\textbf{All database entries are from published empirical
values.} No fitted, computed, or assumed values. The algorithm's
validity rests on the quality of the databases, which are
independently verifiable.

% ============================================================
% PART V: VALIDATED APPLICATIONS
% ============================================================
\newpage
\part{Validated Applications}

\section{Application 1: Room-Temperature Superconductivity}

\subsection{Input}

Target: $T_c \geq 300$~K, ambient pressure, gateable.

\subsection{Forcing Chain}

\begin{enumerate}
  \item A7 (coherence preservation) $\to$ BCS ratio $\leq 5.0$
    $\to$ eliminates all strong-coupling superconductors
  \item A6 (non-trivial winding) $\to$ frustrated geometry
    $\to$ eliminates square, cubic lattices
  \item A8 (local field response) $\to$ $J \geq 100$~meV
    $\to$ eliminates all metals except copper (Cu-O-Cu:
    130~meV)
  \item A6 + A8 $\to$ kagome topology $\to$ demonstrated by
    M\"oller (PCCP 2018) for Cu-O on Au(111)
  \item A5 (dispositional preservation) $\to$ gate doping
    $\to$ eliminates chemical substitution
  \item A3 (boundary compatibility) $\to$ polar oxide
    substrate, lattice match $\to$ LaAlO$_3$(111) at
    5.36~\AA
\end{enumerate}

\subsection{Output}

\textbf{Unique survivor:} Cu$_3$O$_2$ kagome monolayer on
LaAlO$_3$(111), electrostatically gated.

\textbf{Rejected families:} Square cuprates (ratio $>5$),
hydrides (pressure), iron-based ($J$ too low), kagome
pnictides ($J$ too low), twisted bilayer (no high-$J$
mechanism), nickelates (strong coupling AND low $J$).

\textbf{Status:} Patent filed (US Provisional 64/037,691,
April 13, 2026). Blind predictions sealed April 29, 2026:
$T_c = 370 \pm 20$~K. Awaiting experimental synthesis.

% ============================================================
\section{Application 2: Nanoscale Heat Transfer Metamaterials}

\subsection{Input}

Target: Maximum near-field radiative heat transfer enhancement
at 50--1000~nm gaps.

\subsection{Forcing Chain}

\begin{enumerate}
  \item Thermal frequency $\to$ mid-IR operation $\to$
    surface phonon polaritons (SPhP), not plasmons $\to$
    metals eliminated as primary coupling medium
  \item SPhP quality ranking $\to$ SiC forced (Q$\sim$900,
    strongest oscillator strength)
  \item Enhancement beyond planar $\to$ mode confinement
    $\to$ structured metamaterial geometry
  \item Loss minimization $\to$ polar dielectric resonator,
    not metal $\to$ phonon--phonon coupling eliminates Ohmic
    loss $\to$ gold suboptimal
  \item Thermal stability $\to$ all-dielectric survives
    $>$1000$^\circ$C $\to$ gold fails at $\sim$400$^\circ$C
\end{enumerate}

\subsection{Output}

\textbf{Top 3 candidates} (all untested for near-field
radiative heat transfer):
\begin{enumerate}[nosep]
  \item SiC pillar array on SiC substrate (self-matched
    phonon resonance)
  \item hBN nanostructures on SiC (dual-band hyperbolic
    coupling)
  \item SiC phononic crystal slab (slow-light flat bands)
\end{enumerate}

12 total candidates enumerated. All all-phononic.
All fabricable with standard cleanroom processes.

\textbf{Prediction:} All-phononic metamaterials outperform
the plasmonic-phononic hybrid reported in \textit{Nature}
(Wang et al., 2026, 4$\times$ enhancement).

\textbf{Status:} \untested. Published literature supports
each candidate individually. No experimental head-to-head
comparison exists.

% ============================================================
\section{Application 3: Nanoparticle Property Prediction}

\subsection{Method}

The UFCP geometric framework (not CFF filters) applied to
nanoparticle cluster physics. One geometric angle
$\theta = \arctan(1/\varphi)$ (icosahedral frustration, where
$\varphi$ is the golden ratio) generates five projection
operators predicting five branches of physics.

\subsection{Results}

\begin{center}
\begin{tabular}{@{}lrl@{}}
\toprule
\textbf{Property} & \textbf{Predictions} & \textbf{Avg Error} \\
\midrule
Magnetic moments (Fe, Co, Ni, Mn) & 27 & 3.2\% \\
Electron affinity (Au, Ag, Cu) & 23 & 7.1\% \\
Seebeck coefficient (Au$_{13}$) & 1 & 5.8\% \\
HOMO-LUMO gap (Au$_{13}$) & 1 & 0.3\% \\
Seebeck sign (23 elements) & 23 & 0/23 wrong \\
\midrule
\textbf{Total} & \textbf{63+} & \textbf{6.5\%} \\
\bottomrule
\end{tabular}
\end{center}

Zero free parameters. Zero exclusions. 7 elements. 5
geometries. 4 property types. All from one angle.

% ============================================================
\section{Application 4: Phase Transition Detection}

\subsection{Method}

UF-Core kernel (L0--L4) applied to digitized R(T), dielectric,
and DSC data.

\subsection{Results}

\begin{center}
\begin{tabular}{@{}lllc@{}}
\toprule
\textbf{Material} & \textbf{Transition} & \textbf{Precursor} &
  \textbf{Detected?} \\
\midrule
YBCO & $T_c = 93$~K & 13~K lead & \pass \\
MgB$_2$ & $T_c = 39$~K & 18~K lead & \pass \\
CsV$_3$Sb$_5$ & $T_c = 2.5$~K & --- & \pass \\
VO$_2$ & $T_{\text{MIT}} = 340$~K & 20~K lead & \pass \\
BaTiO$_3$ & $T_{\text{Curie}} = 393$~K & 22~K lead & \pass \\
FeSe/STO & $T_c = 65$~K & 20~K lead & \pass \\
Pharma DSC & 5 transitions & --- & 5/5 \pass \\
\bottomrule
\end{tabular}
\end{center}

Same kernel, same parameters, all domains.

% ============================================================
% PART VI: MATHEMATICAL NOTATION AND GLOSSARY
% ============================================================
\newpage
\part{Mathematical Notation and Glossary}

\section{Symbols Used}

\begin{center}
\begin{longtable}{@{}clp{7cm}@{}}
\toprule
\textbf{Symbol} & \textbf{Name} & \textbf{Definition} \\
\midrule
\endhead
$\psi$ & Coherence field & $\psi = \sqrt{\rho}\,e^{i\phi}$;
  complex scalar field on $(M, g_{\mu\nu})$ \\
$\rho$ & Coherent-potential density & $\rho = |\psi|^2 \geq 0$ \\
$\phi$ & Phase & $\phi \in \mathbb{R}$; local phase of
  coherence field \\
$\omega$ & Local tempo & $\omega = \partial_t\phi$; rate of
  phase evolution \\
$S$ & Action & UFCP action functional (Axiom~\ref{ax:action}) \\
$\hat{\Omega}_g$ & Kinetic operator & Hermitian operator on
  $(M, g_{\mu\nu})$; reduces to $-\hbar^2\nabla^2/(2m) + V$
  in non-relativistic limit \\
$U(\rho)$ & Self-interaction & Local potential, bounded below;
  focusing: $U = -\lambda\rho^2/2$ \\
$W$ & Interaction kernel & Symmetric two-point density coupling;
  $W(\mathbf{x}-\mathbf{x}') = W(\mathbf{x}'-\mathbf{x})$ \\
$D_\mu$ & Gauge-covariant derivative &
  $D_\mu = \partial_\mu + i(q/\hbar)A_\mu$ \\
$J^\mu$ & Noether current & Conserved current from $U(1)$
  symmetry: $J^0 = \rho$, $\mathbf{J} = \rho\mathbf{v}$ \\
$\alpha$ & Fine structure constant & $\alpha = 1/N_c^2 \approx
  1/137.036$; from 3D soliton stability \\
$N_c$ & Critical soliton number & $N_c \approx 11.7$; critical
  $L^2$ norm of 3D focusing NLS ground state \\
$\Gamma$ & Geometry factor & $\Gamma = \cos\theta$; ring
  ($\Gamma = 1$), sphere ($\Gamma = 0$) \\
$\lambda$ & Material coupling & Material-specific coupling
  parameter (e.g., electron-phonon coupling) \\
$Q_0$ & Uncoupled observable & Standard value without geometric
  coupling correction \\
$J$ & Superexchange constant & Energy scale of magnetic
  coupling across metal-ligand-metal bridge (meV) \\
$\Delta$ & Superconducting gap & Pairing gap energy \\
$T_c$ & Critical temperature & Phase transition temperature \\
$T_{\text{BKT}}$ & BKT temperature &
  $\pi\hbar^2 n_s/(8m^*k_B)$; 2D superfluid transition \\
$n_s$ & Carrier sheet density & 2D carrier concentration
  (cm$^{-2}$) \\
$m^*$ & Effective mass & Quasiparticle effective mass \\
\midrule
\multicolumn{3}{@{}l}{\textbf{Kernel Symbols}} \\
\midrule
SEV & Structural Event Vector & L0 output:
  $(F, \Delta F, \sigma, \kappa)$ \\
DSF & Decision Surface Field & L4 output:
  $(D_k, M_k, U^*_k, B_k, P_k, C_k, R_{\text{rev},k})$ \\
$D_k$ & Direction & Structural field direction
  ($+1, 0, -1$) \\
$M_k$ & Momentum & Rate of structural change \\
$U^*_k$ & Adjusted uncertainty & Structural confidence;
  peaks at 0.667 at transitions \\
$B_k$ & Breathing & Structural expansion/contraction \\
$P_k$ & Pressure & Transition sensitivity \\
$C_k$ & Complexity & Mosaic divergence \\
$R_{\text{rev},k}$ & Reversal rate & Structural direction
  flip frequency \\
\midrule
\multicolumn{3}{@{}l}{\textbf{CFF Filter Symbols}} \\
\midrule
$r_{max}$ & BCS ratio threshold & Maximum
  $2\Delta/(k_BT_c)$ for target (5.0 for RTSC) \\
$J_{min}$ & Exchange threshold & Minimum superexchange
  for target (100~meV for RTSC) \\
$Q_{\text{SPhP}}$ & Phonon polariton quality &
  $\omega_0/\gamma$ for surface phonon polariton \\
$S_{\text{osc}}$ & Oscillator strength & Relative phonon
  oscillator strength \\
\bottomrule
\end{longtable}
\end{center}

% ============================================================
\section{Acronyms}

\begin{center}
\begin{tabular}{@{}ll@{}}
\toprule
\textbf{Acronym} & \textbf{Expansion} \\
\midrule
CFF & Coherence Field Framework \\
UFCP & Universal Field of Coherent Processes \\
UF-Core & Universal Field Core (L0--L4 kernel) \\
SEV & Structural Event Vector (L0 output) \\
DSF & Decision Surface Field (L4 output) \\
SPhP & Surface Phonon Polariton \\
SPP & Surface Plasmon Polariton \\
HPhP & Hyperbolic Phonon Polariton \\
SES & Secure Envelope System \\
SCE & Structural Cryptographic Engine \\
RTSC & Room-Temperature Superconductivity \\
HTSC & High-Temperature Superconductivity \\
BKT & Berezinskii--Kosterlitz--Thouless \\
BCS & Bardeen--Cooper--Schrieffer \\
NLS & Nonlinear Schr\"odinger (equation) \\
CPT & Charge--Parity--Time (symmetry) \\
NFRHT & Near-Field Radiative Heat Transfer \\
MBE & Molecular Beam Epitaxy \\
RHEED & Reflection High-Energy Electron Diffraction \\
XPS & X-ray Photoelectron Spectroscopy \\
RIE & Reactive Ion Etching \\
FIB & Focused Ion Beam \\
\bottomrule
\end{tabular}
\end{center}

% ============================================================
% PART VII: VERSION HISTORY AND DOCUMENT CONTROL
% ============================================================
\newpage
\part{Document Control}

\section{Version History}

\begin{center}
\begin{tabular}{@{}lp{10cm}@{}}
\toprule
\textbf{Version} & \textbf{Description} \\
\midrule
1.0 (May 2026) & Initial formal specification. Consolidates
  UFCP theoretical foundation (Not-Math v2.0, UFCP Comprehensive
  Spec v1.0), CFF axiom set (A1--A8), UF-Core kernel architecture
  (L0--L4), discovery algorithm (4-stage pipeline), and four
  validated applications (RTSC, metamaterials, nanoparticles,
  phase transitions). Mathematical glossary and acronym table
  included. \\
\bottomrule
\end{tabular}
\end{center}

\section{Source Documents}

This specification consolidates and supersedes:

\begin{enumerate}[nosep]
  \item Not-Math Framework v2.0 (April 2026) --- theoretical
    foundation
  \item UFCP Comprehensive Specification v1.0 (May 2026) ---
    full theoretical treatment with 111 validations
  \item CFF Discovery Algorithm Spec v1.0 (May 2026) ---
    RTSC application and 4-stage pipeline
  \item CFF Metamaterial Heat Transfer Analysis v1.0 (May 2026)
    --- thermal metamaterial application
  \item UF-Core Kernel Spec with SES (Feb 2026) --- L0--L4
    layer definitions and security envelope
  \item DSF-AI Phase Transition Service Spec v1.0 (May 2026) ---
    commercial service architecture
  \item KERNEL\_PHILOSOPHY.md (May 2026) --- operational
    constraints on kernel modification
\end{enumerate}

\section{Intellectual Property}

\begin{itemize}[nosep]
  \item The CFF axiom set (A1--A8) is a \textbf{trade secret}.
  \item The UF-Core kernel (L0--L4) implementation is a
    \textbf{trade secret}.
  \item The DSF 7-tuple field definitions and their structural
    interpretation are \textbf{trade secrets}.
  \item The CFF Discovery Algorithm forcing chain methodology
    is a \textbf{trade secret}.
  \item The UFCP geometric coupling law and its applications
    are \textbf{trade secrets}.
  \item US Provisional Patent 64/037,691 (filed April 13, 2026)
    covers the Cu$_3$O$_2$ kagome superconductor design.
  \item Published results (transition temperatures, precursor
    onsets, regime maps, candidate rankings) are factual
    observations and do not expose trade secrets.
\end{itemize}

\vspace{2em}\hrule\vspace{0.5em}
{\color{nm-gray}\small This document is a trade secret of
Joseph Forrester / DSF-AI. Unauthorized reproduction,
distribution, or disclosure is prohibited. The theoretical
foundation (UFCP) and applied methodology (CFF) are proprietary.
All empirical databases use published values and are
independently verifiable. The kernel implementation and axiom
derivations are not disclosed in any public-facing material.}

\end{document}
